% Options for packages loaded elsewhere
% Options for packages loaded elsewhere
\PassOptionsToPackage{unicode}{hyperref}
\PassOptionsToPackage{hyphens}{url}
\PassOptionsToPackage{dvipsnames,svgnames,x11names}{xcolor}
%
\documentclass[
  11pt,
  a4paper,
  DIV=11,
  numbers=noendperiod]{scrartcl}
\usepackage{xcolor}
\usepackage[top=3cm, bottom=2cm, left=4cm, right=2cm]{geometry}
\usepackage{amsmath,amssymb}
\setcounter{secnumdepth}{-\maxdimen} % remove section numbering
\usepackage{iftex}
\ifPDFTeX
  \usepackage[T1]{fontenc}
  \usepackage[utf8]{inputenc}
  \usepackage{textcomp} % provide euro and other symbols
\else % if luatex or xetex
  \usepackage{unicode-math} % this also loads fontspec
  \defaultfontfeatures{Scale=MatchLowercase}
  \defaultfontfeatures[\rmfamily]{Ligatures=TeX,Scale=1}
\fi
\usepackage{lmodern}
\ifPDFTeX\else
  % xetex/luatex font selection
  \setmainfont[]{Avenir Next}
\fi
% Use upquote if available, for straight quotes in verbatim environments
\IfFileExists{upquote.sty}{\usepackage{upquote}}{}
\IfFileExists{microtype.sty}{% use microtype if available
  \usepackage[]{microtype}
  \UseMicrotypeSet[protrusion]{basicmath} % disable protrusion for tt fonts
}{}
\makeatletter
\@ifundefined{KOMAClassName}{% if non-KOMA class
  \IfFileExists{parskip.sty}{%
    \usepackage{parskip}
  }{% else
    \setlength{\parindent}{0pt}
    \setlength{\parskip}{6pt plus 2pt minus 1pt}}
}{% if KOMA class
  \KOMAoptions{parskip=half}}
\makeatother
% Make \paragraph and \subparagraph free-standing
\makeatletter
\ifx\paragraph\undefined\else
  \let\oldparagraph\paragraph
  \renewcommand{\paragraph}{
    \@ifstar
      \xxxParagraphStar
      \xxxParagraphNoStar
  }
  \newcommand{\xxxParagraphStar}[1]{\oldparagraph*{#1}\mbox{}}
  \newcommand{\xxxParagraphNoStar}[1]{\oldparagraph{#1}\mbox{}}
\fi
\ifx\subparagraph\undefined\else
  \let\oldsubparagraph\subparagraph
  \renewcommand{\subparagraph}{
    \@ifstar
      \xxxSubParagraphStar
      \xxxSubParagraphNoStar
  }
  \newcommand{\xxxSubParagraphStar}[1]{\oldsubparagraph*{#1}\mbox{}}
  \newcommand{\xxxSubParagraphNoStar}[1]{\oldsubparagraph{#1}\mbox{}}
\fi
\makeatother

\usepackage{color}
\usepackage{fancyvrb}
\newcommand{\VerbBar}{|}
\newcommand{\VERB}{\Verb[commandchars=\\\{\}]}
\DefineVerbatimEnvironment{Highlighting}{Verbatim}{commandchars=\\\{\}}
% Add ',fontsize=\small' for more characters per line
\usepackage{framed}
\definecolor{shadecolor}{RGB}{241,243,245}
\newenvironment{Shaded}{\begin{snugshade}}{\end{snugshade}}
\newcommand{\AlertTok}[1]{\textcolor[rgb]{0.68,0.00,0.00}{#1}}
\newcommand{\AnnotationTok}[1]{\textcolor[rgb]{0.37,0.37,0.37}{#1}}
\newcommand{\AttributeTok}[1]{\textcolor[rgb]{0.40,0.45,0.13}{#1}}
\newcommand{\BaseNTok}[1]{\textcolor[rgb]{0.68,0.00,0.00}{#1}}
\newcommand{\BuiltInTok}[1]{\textcolor[rgb]{0.00,0.23,0.31}{#1}}
\newcommand{\CharTok}[1]{\textcolor[rgb]{0.13,0.47,0.30}{#1}}
\newcommand{\CommentTok}[1]{\textcolor[rgb]{0.37,0.37,0.37}{#1}}
\newcommand{\CommentVarTok}[1]{\textcolor[rgb]{0.37,0.37,0.37}{\textit{#1}}}
\newcommand{\ConstantTok}[1]{\textcolor[rgb]{0.56,0.35,0.01}{#1}}
\newcommand{\ControlFlowTok}[1]{\textcolor[rgb]{0.00,0.23,0.31}{\textbf{#1}}}
\newcommand{\DataTypeTok}[1]{\textcolor[rgb]{0.68,0.00,0.00}{#1}}
\newcommand{\DecValTok}[1]{\textcolor[rgb]{0.68,0.00,0.00}{#1}}
\newcommand{\DocumentationTok}[1]{\textcolor[rgb]{0.37,0.37,0.37}{\textit{#1}}}
\newcommand{\ErrorTok}[1]{\textcolor[rgb]{0.68,0.00,0.00}{#1}}
\newcommand{\ExtensionTok}[1]{\textcolor[rgb]{0.00,0.23,0.31}{#1}}
\newcommand{\FloatTok}[1]{\textcolor[rgb]{0.68,0.00,0.00}{#1}}
\newcommand{\FunctionTok}[1]{\textcolor[rgb]{0.28,0.35,0.67}{#1}}
\newcommand{\ImportTok}[1]{\textcolor[rgb]{0.00,0.46,0.62}{#1}}
\newcommand{\InformationTok}[1]{\textcolor[rgb]{0.37,0.37,0.37}{#1}}
\newcommand{\KeywordTok}[1]{\textcolor[rgb]{0.00,0.23,0.31}{\textbf{#1}}}
\newcommand{\NormalTok}[1]{\textcolor[rgb]{0.00,0.23,0.31}{#1}}
\newcommand{\OperatorTok}[1]{\textcolor[rgb]{0.37,0.37,0.37}{#1}}
\newcommand{\OtherTok}[1]{\textcolor[rgb]{0.00,0.23,0.31}{#1}}
\newcommand{\PreprocessorTok}[1]{\textcolor[rgb]{0.68,0.00,0.00}{#1}}
\newcommand{\RegionMarkerTok}[1]{\textcolor[rgb]{0.00,0.23,0.31}{#1}}
\newcommand{\SpecialCharTok}[1]{\textcolor[rgb]{0.37,0.37,0.37}{#1}}
\newcommand{\SpecialStringTok}[1]{\textcolor[rgb]{0.13,0.47,0.30}{#1}}
\newcommand{\StringTok}[1]{\textcolor[rgb]{0.13,0.47,0.30}{#1}}
\newcommand{\VariableTok}[1]{\textcolor[rgb]{0.07,0.07,0.07}{#1}}
\newcommand{\VerbatimStringTok}[1]{\textcolor[rgb]{0.13,0.47,0.30}{#1}}
\newcommand{\WarningTok}[1]{\textcolor[rgb]{0.37,0.37,0.37}{\textit{#1}}}

\usepackage{longtable,booktabs,array}
\usepackage{calc} % for calculating minipage widths
% Correct order of tables after \paragraph or \subparagraph
\usepackage{etoolbox}
\makeatletter
\patchcmd\longtable{\par}{\if@noskipsec\mbox{}\fi\par}{}{}
\makeatother
% Allow footnotes in longtable head/foot
\IfFileExists{footnotehyper.sty}{\usepackage{footnotehyper}}{\usepackage{footnote}}
\makesavenoteenv{longtable}
\usepackage{graphicx}
\makeatletter
\newsavebox\pandoc@box
\newcommand*\pandocbounded[1]{% scales image to fit in text height/width
  \sbox\pandoc@box{#1}%
  \Gscale@div\@tempa{\textheight}{\dimexpr\ht\pandoc@box+\dp\pandoc@box\relax}%
  \Gscale@div\@tempb{\linewidth}{\wd\pandoc@box}%
  \ifdim\@tempb\p@<\@tempa\p@\let\@tempa\@tempb\fi% select the smaller of both
  \ifdim\@tempa\p@<\p@\scalebox{\@tempa}{\usebox\pandoc@box}%
  \else\usebox{\pandoc@box}%
  \fi%
}
% Set default figure placement to htbp
\def\fps@figure{htbp}
\makeatother





\setlength{\emergencystretch}{3em} % prevent overfull lines

\providecommand{\tightlist}{%
  \setlength{\itemsep}{0pt}\setlength{\parskip}{0pt}}



 


\usepackage[document]{ragged2e}
\KOMAoption{captions}{tableheading}
\makeatletter
\@ifpackageloaded{caption}{}{\usepackage{caption}}
\AtBeginDocument{%
\ifdefined\contentsname
  \renewcommand*\contentsname{Table of contents}
\else
  \newcommand\contentsname{Table of contents}
\fi
\ifdefined\listfigurename
  \renewcommand*\listfigurename{List of Figures}
\else
  \newcommand\listfigurename{List of Figures}
\fi
\ifdefined\listtablename
  \renewcommand*\listtablename{List of Tables}
\else
  \newcommand\listtablename{List of Tables}
\fi
\ifdefined\figurename
  \renewcommand*\figurename{Figure}
\else
  \newcommand\figurename{Figure}
\fi
\ifdefined\tablename
  \renewcommand*\tablename{Table}
\else
  \newcommand\tablename{Table}
\fi
}
\@ifpackageloaded{float}{}{\usepackage{float}}
\floatstyle{ruled}
\@ifundefined{c@chapter}{\newfloat{codelisting}{h}{lop}}{\newfloat{codelisting}{h}{lop}[chapter]}
\floatname{codelisting}{Listing}
\newcommand*\listoflistings{\listof{codelisting}{List of Listings}}
\makeatother
\makeatletter
\makeatother
\makeatletter
\@ifpackageloaded{caption}{}{\usepackage{caption}}
\@ifpackageloaded{subcaption}{}{\usepackage{subcaption}}
\makeatother
\usepackage{bookmark}
\IfFileExists{xurl.sty}{\usepackage{xurl}}{} % add URL line breaks if available
\urlstyle{same}
\hypersetup{
  pdftitle={Aufgabenblatt 9: Wequences und Dictionaries},
  colorlinks=true,
  linkcolor={blue},
  filecolor={Maroon},
  citecolor={Blue},
  urlcolor={Blue},
  pdfcreator={LaTeX via pandoc}}


\title{Aufgabenblatt 9: Wequences und Dictionaries}
\author{}
\date{}
\begin{document}
\maketitle


\textbf{Aufgabe 1:} Satzbau

\begin{enumerate}
\def\labelenumi{\alph{enumi})}
\item
  Implementieren Sie ein Funktion namens \texttt{sentence}, welche nach
  Übergabe eines Subjekts, eines Prädikats und eines Objekts aus diesen
  drei Satzbestandteilen einen Satz konstruiert und diesen ausgibt.
\item
  Erstellen Sie mit Hilfe der Funktion ein Python-Progamm, welches die
  drei Satzteile einliest und anschließend ausgibt.
\end{enumerate}

\textbf{Lösung:}

\begin{Shaded}
\begin{Highlighting}[numbers=left,,]
\CommentTok{\# Aufgabe 1}
\end{Highlighting}
\end{Shaded}

\textbf{Aufgabe 2:} C

Gegen sei ein kleiner Supermarkt mit folgenden Produkten. Der Inhaber
des Supermarktes möchte gerne wissen, wie hoch der Verkaufswert aller
seiner Produkte zusammengerechnet beträgt. Schreibe dazu ein
Python-Programm:

\begin{Shaded}
\begin{Highlighting}[numbers=left,,]
\NormalTok{supermarket }\OperatorTok{=}\NormalTok{ \{}\StringTok{"milk"}\NormalTok{: \{}\StringTok{"quantity"}\NormalTok{: }\DecValTok{20}\NormalTok{, }\StringTok{"price"}\NormalTok{: }\FloatTok{1.19}\NormalTok{\},}
               \StringTok{"biscuits"}\NormalTok{:  \{}\StringTok{"quantity"}\NormalTok{: }\DecValTok{32}\NormalTok{, }\StringTok{"price"}\NormalTok{: }\FloatTok{1.45}\NormalTok{\},}
               \StringTok{"butter"}\NormalTok{:  \{}\StringTok{"quantity"}\NormalTok{: }\DecValTok{20}\NormalTok{, }\StringTok{"price"}\NormalTok{: }\FloatTok{2.29}\NormalTok{\},}
               \StringTok{"cheese"}\NormalTok{:  \{}\StringTok{"quantity"}\NormalTok{: }\DecValTok{15}\NormalTok{, }\StringTok{"price"}\NormalTok{: }\FloatTok{1.90}\NormalTok{\},}
               \StringTok{"bread"}\NormalTok{:  \{}\StringTok{"quantity"}\NormalTok{: }\DecValTok{15}\NormalTok{, }\StringTok{"price"}\NormalTok{: }\FloatTok{2.59}\NormalTok{\},}
               \StringTok{"cookies"}\NormalTok{:  \{}\StringTok{"quantity"}\NormalTok{: }\DecValTok{20}\NormalTok{, }\StringTok{"price"}\NormalTok{: }\FloatTok{4.99}\NormalTok{\},}
               \StringTok{"yogurt"}\NormalTok{: \{}\StringTok{"quantity"}\NormalTok{: }\DecValTok{18}\NormalTok{, }\StringTok{"price"}\NormalTok{: }\FloatTok{3.65}\NormalTok{\},}
               \StringTok{"apples"}\NormalTok{:  \{}\StringTok{"quantity"}\NormalTok{: }\DecValTok{35}\NormalTok{, }\StringTok{"price"}\NormalTok{: }\FloatTok{3.15}\NormalTok{\},}
               \StringTok{"oranges"}\NormalTok{:  \{}\StringTok{"quantity"}\NormalTok{: }\DecValTok{40}\NormalTok{, }\StringTok{"price"}\NormalTok{: }\FloatTok{0.99}\NormalTok{\},}
               \StringTok{"bananas"}\NormalTok{: \{}\StringTok{"quantity"}\NormalTok{: }\DecValTok{23}\NormalTok{, }\StringTok{"price"}\NormalTok{: }\FloatTok{1.29}\NormalTok{\}\}\textasciigrave{}\textasciigrave{}\textasciigrave{}}

\OperatorTok{**}\NormalTok{Lösung:}\OperatorTok{**}

\NormalTok{::: \{.cell\}}
\NormalTok{\textasciigrave{}\textasciigrave{}\textasciigrave{} \{.python .cell}\OperatorTok{{-}}\NormalTok{code\}}
\CommentTok{\# Aufgabe 2}
\NormalTok{total\_value }\OperatorTok{=} \DecValTok{0}
\ControlFlowTok{for}\NormalTok{ article, numbers }\KeywordTok{in}\NormalTok{ supermarket.items():}
\NormalTok{    quantity }\OperatorTok{=}\NormalTok{ numbers[}\StringTok{"quantity"}\NormalTok{]}
\NormalTok{    price }\OperatorTok{=}\NormalTok{ numbers[}\StringTok{"price"}\NormalTok{]}
\NormalTok{    product\_price }\OperatorTok{=}\NormalTok{ quantity }\OperatorTok{*}\NormalTok{ price}
\NormalTok{    article }\OperatorTok{=}\NormalTok{ article }\OperatorTok{+} \StringTok{\textquotesingle{}:\textquotesingle{}}
    \BuiltInTok{print}\NormalTok{(}\SpecialStringTok{f"}\SpecialCharTok{\{}\NormalTok{article}\SpecialCharTok{:15s\}}\SpecialStringTok{ }\SpecialCharTok{\{}\NormalTok{product\_price}\SpecialCharTok{:08.2f\}}\SpecialStringTok{"}\NormalTok{)}
\NormalTok{    total\_value }\OperatorTok{+=}\NormalTok{ product\_price}
\BuiltInTok{print}\NormalTok{(}\StringTok{"="}\OperatorTok{*}\DecValTok{24}\NormalTok{)   }
\BuiltInTok{print}\NormalTok{(}\SpecialStringTok{f"Gesamtsumme:    }\SpecialCharTok{\{}\NormalTok{total\_value}\SpecialCharTok{:08.2f\}}\SpecialStringTok{"}\NormalTok{)}
\end{Highlighting}
\end{Shaded}

:::

\textbf{Aufgabe 3:} Wortfrequenzliste mit Python

Erstellen Sie ein Programm freq.py, das aus dem Input eine
Wortfrequenzliste erstellt. Ein Text von einem Stadnard-Input soll
gelesen und in Wörter zerlegt werden. Als Output sollte es allerdings
jedes Wort nur einmal ausgeben, zusammen mit einer Zahl, die angibt, wie
häufig es vorkommt (ähnlich wie sort \textbar{} uniq -c).

Tipp: Verwenden Sie ein Dictionary, um zu zählen, wie oft jedes Wort
vorkommt. Verwenden Sie Wörter als Schlüssel und Zahlen als Werte.

\textbf{Lösung:}

\begin{Shaded}
\begin{Highlighting}[numbers=left,,]
\CommentTok{\# Aufgabe 3}
\end{Highlighting}
\end{Shaded}

\textbf{Aufgabe 4:}




\end{document}
